% =====================================================================
%  Chapter 25 review -- the 14-station revision contract.
%  Revision-only: tables, one-liners, pointers back via \secref/\thmref/
%  \defref/\figref/\coderef/\probref.  No re-derivations.
%  Summit: CF 487E "Tourists" (3200, L6) -- a connected undirected graph with a
%  price w_i per city and two online queries: C a w changes city a's price, and
%  A a b asks for the cheapest souvenir obtainable on SOME simple a-b route.
%  A simple route can detour inside any biconnected block it enters, so the
%  answer is a PATH-MINIMUM on the block-cut tree where each block node stores
%  the min price of its vertices and each vertex node stores its own price.
%  Condense once (Tarjan BCT), then run the chapter's LCT path-min; a price
%  change re-pushes the multiset minima of the (few) blocks containing the city.
%  Both official samples hand-replayed: S1 -> 1, 2 ; S2 -> 2, 1, 5, 3.
%  Exports prob:c23-summit, lst:c23-kit, lst:c23-summit.
%  Mixed set is EXACTLY 10, all live-verified (names/ratings/tags re-pulled from
%  the CF API this session), un-owned, sorted ASCENDING by rating; each hint
%  names the reduction (or the REJECTION) and cites a c23 listing.  CF 1601C is
%  OWNED by Ch.19 -- never linked here.
% =====================================================================

\chaptersection{Chapter 25 Summary}

One sentence: \emph{a splay tree is a binary search tree that shuffles whatever you
just touched to the root, so its \emph{amortised} cost tracks recent use rather than
worst-case depth; a link-cut tree strings those splays along \emph{preferred paths}
and re-preferences them with a single \code{access} loop, which is what lets one
structure reroot, rewire, and fold an aggregate over the current path of a forest
whose edges keep changing (\(\secref{c23-splay}\), \(\secref{c23-makeroot}\)).}
Chapter~22 could only fold a path in a \emph{fixed} tree; this chapter drops that.
Each row below is a pointer, not a proof.

\begin{center}
\renewcommand{\arraystretch}{1.25}
\begin{tabular}{@{}p{0.24\linewidth} p{0.48\linewidth} p{0.20\linewidth}@{}}
\toprule
\textbf{Result} & \textbf{Statement} & \textbf{Where} \\
\midrule
Splay & Repeatedly rotate the accessed key to the root (zig, zig-zig, zig-zag); with
potential \(\Phi=\sum_x\log\size(x)\) every operation is \(\bigO{\log n}\) amortised. &
\defref{c23-splay}, \thmref{c23-splay-amort} \\
Preferred path & The forest is covered by vertex-disjoint \emph{solid} paths, each held in
one aux splay; a non-aux child carries a \emph{path-parent} pointer, so splay-parent and
path-parent are two different roles of \code{fa}. & \defref{c23-prefer}, \figref{c23-aux} \\
Access & \code{access(x)} walks \(x\) to the represented root, at each step splaying and
re-pointing the right child to the last node visited; this makes root\(\to x\) one solid
path and touches \(\bigO{\log n}\) splays amortised. & \defref{c23-access}, \coderef{c23-lct-core} \\
Two-access LCA & \code{access(u)} then climbing \(v\) through path-parents with a splay at
each hop records the last aux root touched -- that node is \(\lca(u,v)\) in a rooted forest. &
\thmref{c23-access-lca}, \coderef{c23-dynalca} \\
Evert / path fold & \code{makeroot(x)} is \code{access(x)} plus a reverse lazy tag on the
whole aux tree; \code{makeroot(u); access(v)} exposes the \(u\)--\(v\) path as \(v\)'s splay,
so any fold monoid reads off \(\mathit{sum}[v]\). & \defref{c23-makeroot}, \thmref{c23-path-query} \\
Subtree fold & A splay sees only \emph{solid} children; the virtual children live in a
separate accumulator \(\mathit{vir}[x]\) that \code{access} updates as it swaps a child
between solid and virtual -- the path fold omits it, the subtree fold includes it. &
\defref{c23-virtual}, \thmref{c23-vsum} \\
Engine map & A link-cut tree earns its constants only when the \emph{topology} moves and a
\emph{path (or subtree) value} is needed online; connectivity goes to ETT/DSU, fixed paths to
HLD, all-pairs to centroid. & \defref{c23-map}, \defref{c23-boundary-pick} \\
\bottomrule
\end{tabular}
\end{center}

\medskip
{\small\textbf{Algorithms from this chapter.}}
\begin{center}
\renewcommand{\arraystretch}{1.2}
\begin{tabular}{@{}p{0.36\linewidth} l l >{\raggedright\arraybackslash}p{0.20\linewidth}@{}}
\toprule
\textbf{Algorithm} & \textbf{Prep} & \textbf{Per op} & \textbf{Code} \\
\midrule
Bottom-up splay (rotate to root) & -- & \(\bigO{\log n}\) amort & \coderef{c23-splay-rot} \\
Top-down splay (no recursion/parents) & -- & \(\bigO{\log n}\) amort & \coderef{c23-topdown} \\
SPOJ ORDERSET (order-statistics set) & \(\bigO{1}\) & \(\bigO{\log n}\) amort & \coderef{c23-orderset} \\
LCT core: access / findroot (rooted) & \(\bigO{n}\) & \(\bigO{\log n}\) amort & \coderef{c23-lct-core} \\
SPOJ DYNALCA (link/cut/lca) & \(\bigO{n}\) & \(\bigO{\log n}\) amort & \coderef{c23-dynalca} \\
Evert + path-sum LCT & \(\bigO{n}\) & \(\bigO{\log n}\) amort & \coderef{c23-evert}, \coderef{c23-pathsum} \\
Subtree-sum LCT (\(\mathit{vir}\) sibling) & \(\bigO{n}\) & \(\bigO{\log n}\) amort & \coderef{c23-subtree} \\
Engine pick table (is it even LCT?) & -- & -- & \defref{c23-boundary-pick}, \defref{c23-map} \\
Link-cut kit (one listing) & -- & -- & \coderef{c23-kit} \\
CF 487E (Summit) -- BCT + LCT path-min & \(\bigO{n+m}\) & \(\bigO{\log n}\) amort & \coderef{c23-summit} \\
\bottomrule
\end{tabular}
\end{center}

\chaptersection{Master Decision Flowchart}

Post-study ladder: walk top to bottom, stop at the first branch that matches. It decides
\emph{whether the topology truly moves} (the LCT's only excuse for existing) and, if so,
\emph{which operation} -- \code{access} alone, \code{makeroot} plus a path fold, or the
\(\mathit{vir}\) subtree sibling -- the query needs. Everything that fails the top gate is
handed back to a static engine (HLD \(\chref{c16}\), Euler+Fenwick \(\chref{c24}\), offline DSU
\(\chref{c28}\)). A \emph{graph} that is not a forest must be \emph{condensed} first -- the Summit.

\begin{lstlisting}[language={}, basicstyle=\ttfamily\small, numbers=none, frame=single]
  Does the structure answer queries about a FOREST WHOSE EDGES CHANGE (link / cut)?
  |
  |-- The tree is FIXED -- no edge is ever added or deleted, only VALUES move?
  |     +-- YES: NOT link-cut (def:c23-boundary-pick).  Path range add/sum -> HLD
  |         (Ch.22, c16); subtree range -> Euler tour + Fenwick (Ch.5, c24).
  |         Reaching for an LCT here pays splay constants for nothing.
  |
  |-- Only CONNECTIVITY (are u,v in one tree?) -- never a path value?
  |     +-- offline (all edits known first): rollback DSU / segment-tree-on-time
  |         (Ch.17 c28, Ch.24 c18).  Fastest, least code.  NOT LCT.
  |     +-- online: Euler-tour tree (Ch.5, c24) or an LCT findroot -- but if no
  |         path value is needed, the ETT is the smaller, safer structure.
  |
  |-- A PATH (or SUBTREE) VALUE on a forest whose edges DO change, online: THIS chapter.
  |     First own the splay engine (25.1): rot / splay, zig-zig vs zig-zag, size
  |     fields for k-th.  Without the amortised argument (thm:c23-splay-amort) the
  |     LCT is just a slow BST.  lst:c23-splay-rot, lst:c23-topdown.
  |
  |     |-- ROOTED growing forest, need LCA / root / connected?  (25.2)
  |     |     access(x) exposes root->x; findroot = access then leftmost then splay;
  |     |     lca(u,v) = access(u) then the two-access climb on v.  lst:c23-lct-core,
  |     |     lst:c23-dynalca.  NO reverse tag here -- keep the root fixed.
  |     |
  |     |-- ARBITRARY path u~v fold (sum/min/max) with rerooting?  (25.3)
  |     |     makeroot(x)=access+apply-reverse; query = makeroot(u); access(v);
  |     |     read sum[v].  link = makeroot(u); fa[u]=v.  cut = makeroot(u);
  |     |     access(v); detach.  Push rev tags top-down BEFORE any rotate.
  |     |     lst:c23-evert, lst:c23-pathsum.
  |     |
  |     |-- SUBTREE aggregate on a moving tree?  (25.4)
  |     |     add the vir[] sibling: access moves a child between solid and virtual,
  |     |     so vir[x]-=sub[child-in], vir[x]+=sub[old-right], pull after.  sub[]
  |     |     folds val+vir+solid; the PATH fold sum[] must NOT include vir.
  |     |     lst:c23-subtree.  (thm:c23-vsum)
  |     |
  |     |-- The input is a GRAPH, not a forest?  Condense, THEN run the LCT.
  |         Block-cut (biconnected) tree or a maintained spanning forest turns the
  |         query into a path fold on a tree -- this is exactly the Summit, CF 487E:
  |         "cheapest city on SOME simple route" -> path-min on the block-cut tree
  |         where a block node carries the min price of its vertices.  lst:c23-summit.
  |
  |-- NONE of the above -- you wanted a static engine.  See 25.5 (def:c23-boundary-pick):
        fixed tree -> HLD / Euler; counting over ALL paths -> centroid (Ch.24, c18);
        per-node state -> tree DP (Ch.13, c10); offline node-sets -> Mo (Ch.27, c21).
\end{lstlisting}

\chaptersection{Common Mistakes Checklist}

\begin{itemize}[nosep, leftmargin=1.4em]
  \item \textbf{Confusing the splay-parent with the path-parent.} \code{fa[x]} is a splay edge
        when \(x\) is an aux child and a \emph{path-parent} when \(x\) is an aux root; the test
        is \(\mathit{isRoot}(x)\): \code{fa}==0 \emph{or} \code{fa} has no child equal to \(x\).
        Writing \(\mathit{isRoot}(x)\equiv(\mathit{fa}[x]=0)\) breaks every access. \(\defref{c23-prefer}\), \(\coderef{c23-lct-core}\).
  \item \textbf{Not pushing the reverse tag top-down before rotating.} After \code{makeroot} the
        whole aux subtree is lazily flipped; a \code{splay} that rotates before pushing \code{rev}
        to the ancestors reattaches mirrored children -- \textbf{WA} that only shows on deep
        paths. Push from the aux root down before the zig-zig loop. \(\defref{c23-makeroot}\), \(\coderef{c23-evert}\).
  \item \textbf{Forgetting \code{pull} after \code{access} re-points the right child.} The line
        \(\mathit{ch}[x][1]\leftarrow\mathit{last}\) changes the splay subtree, so \(\mathit{sum}\)/\(\mathit{mn}\)
        must be recomputed before the next hop, or the aggregate is stale. \(\defref{c23-access}\), \(\coderef{c23-lct-core}\).
  \item \textbf{\code{link} without \code{makeroot}.} \code{link(u,v)} must expose \(u\) as a
        represented root first (\code{makeroot(u); fa[u]=v}); linking a node that already has a
        parent silently welds two trees into a structure that is not a forest. \(\coderef{c23-evert}\).
  \item \textbf{\code{findroot} not splaying at the end.} Walking to the leftmost node then
        \emph{not} splaying it destroys the amortised bound and leaves the exposed root buried;
        always \code{splay} the returned root. \(\coderef{c23-lct-core}\).
  \item \textbf{Reading the subtree fold off the path aggregate (or vice-versa).} \(\mathit{sum}\)
        omits \(\mathit{vir}\); \(\mathit{sub}\) includes it. A subtree query answered from
        \(\mathit{sum}\), or a path query from \(\mathit{sub}\), mis-counts every branch. \(\thmref{c23-vsum}\), \(\coderef{c23-subtree}\).
  \item \textbf{Wrong \(\mathit{vir}\) update order in \code{access}.} Compute
        \(\mathit{sub}[\mathit{last}]\) and \(\mathit{sub}[r]\) \emph{before} re-pointing
        \(\mathit{ch}[x][1]\); after the re-point the demoted child's value is unreachable, so a
        late read returns garbage. \(\thmref{c23-vsum}\), \(\coderef{c23-subtree}\).
  \item \textbf{Using an LCT on a tree that never changes.} If no edge is ever added or deleted,
        HLD has the same asymptotics with a debuggable array segment tree; the LCT is pure
        overhead and a larger bug surface. \(\defref{c23-boundary-pick}\), \(\secref{c23-boundary}\).
  \item \textbf{Sentinel node \(0\) not neutral.} Node \(0\) is the null: its \(\mathit{size}=0\)
        for order statistics and its \(\mathit{mn}=+\infty\) for a path minimum; forget the
        neutral value and the first rotation corrupts the aggregate. \(\coderef{c23-orderset}\), \(\coderef{c23-summit}\).
\end{itemize}

\chaptersection{WA Clinic}

Verdicts are symptoms; the causes are the short list this chapter drilled. Diagnose top-down --
every leaf names a listing.

\begin{figure}[H]
\centering
\begin{tikzpicture}[diagram,
    v/.style={boxnode, fill=primaryblue!12, draw=primaryblue, text width=22mm, align=center},
    b/.style={boxnode, fill=difforange!12, draw=difforange, text width=12mm},
    l/.style={boxnode, fill=successgreen!10, draw=successgreen, text width=48mm,
              font=\scriptsize, align=left}]
  \node[v] (verdict) {\textbf{Verdict on the judge}};
  \node[b, below left=8mm and 5mm of verdict] (wa) {\textbf{WA}};
  \node[b, below right=8mm and 5mm of verdict] (tle) {\textbf{TLE / RE}};
  \node[l, below right=1mm and 2mm of wa.south] (wa1)
    {the reverse tag was never pushed before a rotate -- deep paths come back mirrored, so a
     path fold reads the wrong vertices (\(\defref{c23-makeroot}\), \(\coderef{c23-evert}\))};
  \node[l, below left=1mm and 2mm of wa.south] (wa2)
    {\(\mathit{isRoot}\) written as \code{fa==0}, conflating path-parent with splay-parent, or a
     subtree query answered from the path fold (\(\coderef{c23-lct-core}\), \(\thmref{c23-vsum}\))};
  \node[l, below right=1mm and 2mm of tle.south] (tle1)
    {no amortised discipline: \code{findroot} fails to splay, or the tree never changes so an LCT
     is paying splay constants where HLD was free (\(\coderef{c23-lct-core}\), \(\defref{c23-boundary-pick}\))};
  \node[l, below left=1mm and 2mm of tle.south] (tle2)
    {recursion depth on a chain splay blows the stack, or the sentinel \(0\) carries a non-neutral
     aggregate (\(\coderef{c23-topdown}\), \(\coderef{c23-summit}\))};
  \draw[flow] (verdict) -- (wa);
  \draw[flow] (verdict) -- (tle);
  \draw[flow] (wa) -- (wa1);
  \draw[flow] (wa) -- (wa2);
  \draw[flow] (tle) -- (tle1);
  \draw[flow] (tle) -- (tle2);
\end{tikzpicture}
\caption*{WA: the lazy reverse tag is not pushed before rotating, or \(\mathit{isRoot}\) mistakes a
path-parent for a splay-parent, or the path/subtree folds are swapped. TLE or RE: \code{findroot}
forgets its final splay, the topology never actually changes (use HLD), or the null sentinel is not
neutral so the first rotate corrupts an aggregate.}
\end{figure}

\chaptersection{Quick Reference Card}

\begin{quickref}[Chapter 25 at a glance]
\textbf{A link-cut tree is a set of splay trees, one per preferred path of a forest whose edges
change; \code{access} re-preferences the root-to-node path in one loop, \code{makeroot} reroots with
a lazy flip, and any path (or, with \(\mathit{vir}\), subtree) fold answers in \(\bigO{\log n}\) amortised. }

\medskip
\textbf{Splay (25.1).} Rotate the touched key to the aux root: zig (root child), zig-zig (two same-side,
one rotate of the parent first), zig-zag (opposite sides, two rotates of \(x\)). With
\(\Phi=\sum_x\log\size(x)\) this is \(\bigO{\log n}\) amortised; keep \(\mathit{size}\) for k-th.
(\(\defref{c23-splay}\), \(\thmref{c23-splay-amort}\), \(\coderef{c23-splay-rot}\), \(\coderef{c23-topdown}\)).

\medskip
\textbf{Access / rooted (25.2).} \(\mathit{isRoot}(x)\) is a splay test, not \code{fa==0}. \code{access(x)}:
splay, set right child to \code{last}, pull, step to \code{fa}. \code{findroot}: access then leftmost then
splay. \code{lca(u,v)}: access(u); climb \(v\) by path-parent splaying, return last.
(\(\defref{c23-access}\), \(\thmref{c23-access-lca}\), \(\coderef{c23-lct-core}\), \(\coderef{c23-dynalca}\)).

\medskip
\textbf{Evert / paths (25.3).} \code{makeroot(x)}=access+flip (swap children, toggle rev). \code{query(u,v)}=
makeroot(u); access(v); read \(\mathit{fold}[v]\). \code{link}: makeroot(u); fa[u]=v. \code{cut}: makeroot(u);
access(v); detach \(v\)'s left child. Push rev top-down before rotating.
(\(\defref{c23-makeroot}\), \(\thmref{c23-path-query}\), \(\coderef{c23-evert}\), \(\coderef{c23-pathsum}\)).

\medskip
\textbf{Subtree / boundary (25.4).} Keep \(\mathit{vir}[x]\) = fold over virtual children; \code{access}
moves a child between the solid splay and \(\mathit{vir}\) (subtract the promoted, add the demoted).
\(\mathit{sub}=\mathit{val}+\mathit{vir}+\) solid, \(\mathit{sum}\) omits \(\mathit{vir}\).
(\(\defref{c23-virtual}\), \(\thmref{c23-vsum}\), \(\coderef{c23-subtree}\)).

\medskip
\textbf{The boundary (25.5).} Fixed tree \(\to\) HLD (\(\chref{c16}\)) / Euler+Fenwick (\(\chref{c24}\));
connectivity only, offline \(\to\) rollback DSU (\(\chref{c28}\)); connectivity only, online \(\to\)
Euler-tour tree; all-pairs counting \(\to\) centroid (\(\chref{c18}\)); \emph{moving forest + online path
value} \(\to\) LCT (\(\defref{c23-boundary-pick}\), \(\defref{c23-map}\)).
\end{quickref}

\begin{lstlisting}[caption={Chapter 25 kit: the splay rotate-and-splay, the rooted LCT core (isRoot, access, findroot), evert with its reverse tag plus link/cut/path-query, the subtree fold with its vir sibling, the two-access LCA, and a path-min variant -- one compact skeleton each, in one place. ASCII only, identifiers spelled without underscores.}, label={lst:c23-kit}]
// ===== SPLAY (25.1) =====
int ch[MAXN][2], fa[MAXN];
bool isRoot(int x){ int f=fa[x]; return f==0 || (ch[f][0]!=x && ch[f][1]!=x); }
void rot(int x){ int y=fa[x], z=fa[y], d=(ch[y][1]==x);
    if(!isRoot(y)) ch[z][ch[z][1]==y]=x; fa[x]=z;
    ch[y][d]=ch[x][d^1]; if(ch[x][d^1]) fa[ch[x][d^1]]=y;
    ch[x][d^1]=y; fa[y]=x; pull(y); pull(x); }
void splay(int x){ pushPath(x);
    while(!isRoot(x)){ int y=fa[x], z=fa[y];
        if(!isRoot(y)) rot( (ch[y][0]==x)==(ch[z][0]==y) ? y : x );  // zig-zig vs zig-zag
        rot(x); } }

// ===== ROOTED LCT CORE (25.2) : NO reverse tag, root is fixed =====
void access(int x){ for(int last=0; x; last=x, x=fa[x]){ splay(x); ch[x][1]=last; pull(x);} }
int findroot(int x){ access(x); while(ch[x][0]) { push(x); x=ch[x][0]; } splay(x); return x; }

// ===== EVERT + PATH FOLD (25.3) : sum is a generic monoid fold =====
bool rev[MAXN];
void pushRev(int x){ if(!x) return; swap(ch[x][0],ch[x][1]); rev[x]^=1; }
void push(int x){ if(rev[x]){ pushRev(ch[x][0]); pushRev(ch[x][1]); rev[x]=0; } }
void makeroot(int x){ access(x); pushRev(x); }
void link(int u,int v){ makeroot(u); fa[u]=v; }
void cut(int u,int v){ makeroot(u); access(v); if(ch[v][0]==u && !ch[u][1]){ ch[v][0]=0; fa[u]=0; pull(v);} }
ll query(int u,int v){ makeroot(u); access(v); return sum[v]; }   // fold over u~v

// ===== SUBTREE FOLD with vir sibling (25.4) =====
void accessV(int x){ for(int last=0; x; last=x, x=fa[x]){ splay(x);
    int r=ch[x][1];                                        // old solid child -> virtual
    if(r) vir[x]+=sub[r]; if(last) vir[x]-=sub[last];
    ch[x][1]=last; pull(x); } }                            // pull: sum=val+solid ; sub=val+vir+solid
ll subtree(int x){ accessV(x); return sub[x]; }

// ===== TWO-ACCESS LCA (25.2, rooted) =====
int lca(int u,int v){ if(findroot(u)!=findroot(v)) return 0;
    access(u); int last=0; for(int x=v; x; x=fa[x]){ splay(x); last=x; } splay(last); return last; }

// ===== PATH-MIN variant (Summit 487E) : mn[0]=+INF, mn=min(val, children) =====
void pullMin(int x){ mn[x]=val[x]; if(ch[x][0]) mn[x]=std::min(mn[x],mn[ch[x][0]]);
    if(ch[x][1]) mn[x]=std::min(mn[x],mn[ch[x][1]]); }
ll pathMin(int u,int v){ makeroot(u); access(v); return mn[v]; }
\end{lstlisting}

\chaptersection{Write-From-Memory Drills}

Five signature skeletons, each lifted from a listing already in the chapter. Cover the marked line,
reproduce from memory, then check the answer one line below.

\begin{lstlisting}[caption={Drill 1 -- splay step: decide zig-zig vs zig-zag, then rotate (from lst:c23-splay-rot)}]
void splay(int x){ pushPath(x);
    while(!isRoot(x)){
        int y = fa[x], z = fa[y];
        if(!isRoot(y))                       // grandparent exists -> two-level case
            rot( (ch[y][0]==x) == (ch[z][0]==y) ? y : x );   // same side: rotate y; else rotate x
        rot(x);                              // always finish by rotating x
    }
}
\end{lstlisting}
\noindent{\scriptsize Key: the parity test \((\mathit{ch}[y][0]{=}x)==(\mathit{ch}[z][0]{=}y)\) is
zig-zig (grandparent and parent go the same way, rotate the \emph{parent} first) versus zig-zag
(opposite ways, rotate \emph{x} first); push the lazy tags from the aux root down before the loop
(\(\defref{c23-splay}\), \(\thmref{c23-splay-amort}\), \(\coderef{c23-splay-rot}\)).}

\begin{lstlisting}[caption={Drill 2 -- the access loop: the whole engine in four lines (from lst:c23-lct-core)}]
void access(int x){
    for(int last = 0; x; last = x, x = fa[x]){
        splay(x);            // x to the top of its aux tree
        ch[x][1] = last;     // re-prefer: the path below x becomes the right spine
        pull(x);             // refresh x's aggregate after the child swap
    }
}
\end{lstlisting}
\noindent{\scriptsize Key: after the step \(x\leftarrow\mathit{fa}[x]\), the old right child is
\emph{demoted} (it becomes virtual) and \code{last} is \emph{promoted} to solid; that single
re-point, \(\bigO{\log n}\) times, is what exposes root\(\to\)the original \(x\) as one solid path
(\(\defref{c23-access}\), \(\coderef{c23-lct-core}\)).}

\begin{lstlisting}[caption={Drill 3 -- evert and read a path fold (from lst:c23-evert)}]
void makeroot(int x){ access(x); pushRev(x); }        // x becomes the represented root
ll  query(int u,int v){ makeroot(u); access(v); return sum[v]; }  // u~v is v's aux subtree
void link(int u,int v){ makeroot(u); fa[u]=v; }       // hang u (now a root) under v
void cut(int u,int v){ makeroot(u); access(v);
    if(ch[v][0]==u && !ch[u][1]){ ch[v][0]=0; fa[u]=0; pull(v); } }
\end{lstlisting}
\noindent{\scriptsize Key: \code{makeroot} is exactly \code{access} plus a lazy reversal of the whole
aux tree; once \(u\) is the root, \code{access(v)} makes the \(u\)--\(v\) route the right spine, so
the fold sits at \(v\). The \code{cut} guard (u is v's left child with no right child) prevents
severing a non-edge (\(\defref{c23-makeroot}\), \(\thmref{c23-path-query}\), \(\coderef{c23-evert}\)).}

\begin{lstlisting}[caption={Drill 4 -- the two-access LCA on a rooted forest (from lst:c23-dynalca)}]
int lca(int u, int v){
    if(findroot(u) != findroot(v)) return 0;          // different trees -> no LCA
    access(u);                                        // expose root->u as one solid path
    int last = 0;
    for(int x = v; x; x = fa[x]) { splay(x); last = x; }  // climb v's path-parents
    splay(last);                                      // keep the amortised bound
    return last;                                      // first node already in u's path = LCA
}
\end{lstlisting}
\noindent{\scriptsize Key: after \code{access(u)} the nodes root\(\to u\) sit on solid paths; climbing
\(v\) by \code{fa} and splaying each hop, the \emph{last} aux root reached before you fall off to the
shared part is \(\lca(u,v)\) -- the two-access method, no reverse tag, root fixed
(\(\thmref{c23-access-lca}\), \(\coderef{c23-dynalca}\)).}

\begin{lstlisting}[caption={Drill 5 -- subtree fold via the vir sibling in access (from lst:c23-subtree)}]
void accessSub(int x){
    for(int last = 0; x; last = x, x = fa[x]){
        splay(x);
        int r = ch[x][1];                 // solid child about to be demoted to virtual
        if(r)    vir[x] += sub[r];        // add its WHOLE subtree into the virtual pool
        if(last) vir[x] -= sub[last];     // remove the promoted child from the pool
        ch[x][1] = last;
        pull(x);                          // sum=val+solid (no vir) ; sub=val+vir+solid
    }
}
\end{lstlisting}
\noindent{\scriptsize Key: the path fold \(\mathit{sum}\) ignores \(\mathit{vir}\) because off-path
branches are not on a path; the subtree fold \(\mathit{sub}\) must add \(\mathit{vir}\) because those
branches \emph{are} in the subtree. Read \(\mathit{sub}[\mathit{last}]\) and \(\mathit{sub}[r]\) before
re-pointing \(\mathit{ch}[x][1]\) (\(\thmref{c23-vsum}\), \(\coderef{c23-subtree}\)).}


\chaptersection{Complete Problem Checklist}

Every problem this chapter offers -- two fully solved in-section, worked traces on the running
forest \(R\), and one Summit. Tick \emph{only} after an unaided solve; a hinted attempt counts as half.

\begin{center}
\footnotesize
\setlength{\tabcolsep}{4pt}
\renewcommand{\arraystretch}{1.25}
\begin{tabular}{@{}c p{0.16\linewidth} l p{0.20\linewidth} >{\raggedright\arraybackslash}p{0.26\linewidth} c@{}}
\toprule
\textbf{\#} & \textbf{Problem} & \textbf{Source} & \textbf{Key technique} & \textbf{What to look for} & \textbf{Done?} \\
\midrule
\multicolumn{6}{l}{\textbf{Section 25.1 -- the splay engine}} \\
W25.1 & replay splay on keys \(\{1,2,3,4\}\) & worked & \lv{4} zig / zig-zig / zig-zag & access order reshapes depth; \(\mathit{size}\) kept on every rotate; \(\Phi\) drops & \(\square\) \\
W25.1b & Order Statistic Set & SPOJ & \lv{4} splay + \(\mathit{size}\) kth & dedup insert, kth by size, count below \(x\), erase; sample \(\to\) 1,2,2,invalid & \(\square\) \\
\midrule
\multicolumn{6}{l}{\textbf{Section 25.2 -- preferred paths and access}} \\
W25.2 & \code{access(7)} on forest \(R\) & worked & \lv{4} re-preference loop & root\(\to\)7 solid; siblings demoted to virtual; one aux splay per path & \(\square\) \\
W25.2b & Dynamic LCA & SPOJ & \lv{5} two-access LCA & link/cut/\emph{lca}, rooted, no reverse tag; sample \(\to\) 1,2,3,2 & \(\square\) \\
\midrule
\multicolumn{6}{l}{\textbf{Section 25.3 -- evert, link, cut, path fold (worked only)}} \\
W25.3 & path sum on chain 1--2--3--4 & worked & \lv{5} makeroot + fold & query(1,4)=10; add node \(\to\) 15; cut splits \(\{1,2\},\{3,4\}\); push rev before rotate & \(\square\) \\
\midrule
\multicolumn{6}{l}{\textbf{Section 25.4 -- virtual children and connectivity (worked only)}} \\
W25.4 & subtree sum after \code{access} & worked & \lv{5} \(\mathit{vir}\) sibling & demote solid \(\Rightarrow\) add to \(\mathit{vir}\); promote \(\Rightarrow\) subtract; \(\mathit{sub}\ne\mathit{sum}\) & \(\square\) \\
\midrule
\multicolumn{6}{l}{\textbf{Section 25.5 -- the engine boundary (worked only)}} \\
W25.5 & route six specs off the LCT & worked & \lv{5} which engine & fixed \(\to\) HLD/Euler; offline conn. \(\to\) DSU; all-pairs \(\to\) centroid; moving path \(\to\) LCT & \(\square\) \\
\midrule
\multicolumn{6}{l}{\textbf{Summit (chapter end)}} \\
W25 & Tourists & CF 487E & \lv{6} block-cut tree + LCT path-min & condense graph to BCT, block node \(=\) min price; reruns \(\to\) S1: 1,2; S2: 2,1,5,3 & \(\square\) \\
\bottomrule
\end{tabular}
\end{center}

\chaptersection{Recommended Study Order}

Five focused days. Day~1 internalises the splay rotations and ORDERSET; day~2 the access loop and
DYNALCA; day~3 evert and the path fold; day~4 the \(\mathit{vir}\) subtree sibling and the engine
boundary; day~5 the Summit and the mixed set.

\begin{center}
\small
\renewcommand{\arraystretch}{1.25}
\begin{tabular}{@{}llp{0.55\linewidth}@{}}
\toprule
\textbf{Day} & \textbf{Focus} & \textbf{Tasks} \\
\midrule
1 & Read + do 25.1 & hand-splay keys \(\{1,2,3,4\}\) after accessing 1 then 4, name each zig/zig-zig/zig-zag (\(\figref{c23-splaycases}\)); write \(\coderef{c23-splay-rot}\) cold; then SPOJ ORDERSET with the \(\mathit{size}\) kth (\(\coderef{c23-orderset}\)) \\
2 & Read + do 25.2 & trace \code{access(7)} on \(R\), marking which children become virtual (\(\coderef{c23-lct-core}\)); then SPOJ DYNALCA, replaying the sample to \(1,2,3,2\) (\(\coderef{c23-dynalca}\)) \\
3 & Read + do 25.3 & trace \code{makeroot}/\code{query} on the chain 1--2--3--4 to \(10\) then \(15\) (\(\coderef{c23-evert}\), \(\coderef{c23-pathsum}\)); write the reverse-tag pushdown cold \\
4 & Read + do 25.4 + 25.5 & add the \(\mathit{vir}\) sibling and re-trace the subtree query (\(\coderef{c23-subtree}\)); place the six specs on the pick (\(\defref{c23-boundary-pick}\)) \\
5 & Summit + mixed & condense the CF 487E sample graph to its block-cut tree and replay the answers \(2,1,5,3\) (\(\probref{c23-summit}\)); then the first 3 mixed entries \\
\bottomrule
\end{tabular}
\end{center}

\chaptersection{Difficulty Progression}

Tiers are competency gates -- each row names the \emph{skill} the tier unlocks.

\begin{center}
\small
\renewcommand{\arraystretch}{1.3}
\begin{tabular}{@{}lp{0.32\linewidth}p{0.40\linewidth}@{}}
\toprule
\textbf{Tier} & \textbf{Problems} & \textbf{What passing this tier proves} \\
\midrule
\lv{3} \ \(\bigstar\bigstar\bigstar\)
  & a single splay rotate; one \code{access} step
  & you keep the splay/path-parent distinction and refresh the aggregate after each child swap \\
\lv{4} \ \(\bigstar\bigstar\bigstar\bigstar\)
  & ORDERSET; \code{access}/\code{findroot} on \(R\)
  & you maintain \(\mathit{size}\) for k-th and expose a root-to-node path correctly \\
\lv{5} \ \(\bigstar\bigstar\bigstar\bigstar\bigstar\)
  & DYNALCA, evert path fold, \(\mathit{vir}\) subtree, mixed at CF~2100--2400
  & you push reverse tags before rotating and keep \(\mathit{sub}\) distinct from \(\mathit{sum}\) \\
\lv{6} \ \(\bigstar\bigstar\bigstar\bigstar\bigstar\bigstar\)
  & CF 487E (Summit), the CF~3000--3400 mixed set
  & you \emph{condense a graph to a block-cut tree}, then run an LCT path fold on it \\
\bottomrule
\end{tabular}
\end{center}

\chaptersection{Reduction Gallery}

Every entry is ``one dynamic-forest shape \(\to\) one LCT move (or an honest hand-off).'' Each line
collapses a problem family onto a listing already in the chapter.

\begin{itemize}[nosep, leftmargin=1.4em]
  \item ``maintain an ordered set with k-th / rank / predecessor under insert--delete'' (SPOJ ORDERSET)
        \(\to\) a splay keyed by value with a \(\mathit{size}\) field; k-th walks by subtree sizes
        (\(\defref{c23-splay}\), \(\coderef{c23-orderset}\)).
  \item ``link/cut a \emph{rooted} forest, ask LCA / connectivity online'' (SPOJ DYNALCA) \(\to\) the
        rooted core: \code{access} + the two-access climb, \code{findroot} for the same-tree test
        (\(\thmref{c23-access-lca}\), \(\coderef{c23-dynalca}\)).
  \item ``link/cut online, fold sum/min/max over the current path \(u\)--\(v\)'' \(\to\)
        \code{makeroot(u); access(v); read fold[v]}, with a lazy reverse tag (\(\thmref{c23-path-query}\),
        \(\coderef{c23-evert}\)).
  \item ``subtree aggregate on a moving tree'' \(\to\) the \(\mathit{vir}\) sibling kept current in
        \code{access}; subtree fold \(\mathit{sub}\) includes it, path fold \(\mathit{sum}\) does not
        (\(\thmref{c23-vsum}\), \(\coderef{c23-subtree}\)).
  \item ``cheapest vertex on \emph{some} simple route \(a\to b\) in a graph, with price changes''
        (CF 487E, Summit) \(\to\) condense to the block-cut tree, a block node carrying its
        vertices' min, then an LCT path-min (\(\coderef{c23-summit}\)).
  \item ``connectivity only, edits known offline'' \(\to\) \emph{not} an LCT: rollback DSU /
        segment-tree-on-time (\(\defref{c23-map}\), \(\secref{c23-boundary}\)).
  \item ``path/subtree ranges on a \emph{fixed} tree'' \(\to\) \emph{not} dynamic: HLD (\(\chref{c16}\))
        or Euler+Fenwick (\(\chref{c24}\)) (\(\defref{c23-boundary-pick}\)).
  \item ``counting / aggregating over \emph{all} simple paths of a fixed tree'' \(\to\) centroid
        decomposition (\(\chref{c18}\)), not a per-path fold (\(\secref{c23-boundary}\)).
\end{itemize}

\chaptersection{Mixed Problem Set}

Ten entries, each live-verified against the current Codeforces problem set (names, ratings, and tags
re-pulled this session) and absent from every other roster in the book (the owned IDs -- CF 1601C in
Ch.~19, and this chapter's own SPOJ solveds -- are excluded). None says ``use a link-cut tree'';
several deliberately \emph{do not} want one, and the drill is to say so. Each hint names the
reduction (or the rejection) and the Chapter~25 listing it leans on, ordered ascending by rating.

\begin{mixedset}
  \practice{Chain Queries}{CF 1975E}
    {https://codeforces.com/problemset/problem/1975/E}{CF 2100 L5}
    {A \emph{fixed} tree with queries resolved by a DFS order plus a binary search on the answer --
     no edge is ever linked or cut. Run the top gate of the pick (\(\defref{c23-boundary-pick}\)):
     the topology is static, so this is a static path/chain job for HLD (\(\chref{c16}\)) or an
     Euler tour, \emph{not} a link-cut tree. Use this entry to practise REJECTING the LCT for want
     of a moving edge; contrast \(\secref{c23-boundary}\)'s first spec.}
  \practice{Avoid Rainbow Cycles}{CF 1408E}
    {https://codeforces.com/problemset/problem/1408/E}{CF 2400 L5}
    {Greedy over a graph whose vertices and colours form a bipartite structure, unioning edges with a
     DSU so as to avoid ``rainbow cycles'' -- one maximum-weight forest pass, no online path value.
     The pick hands this to DSU + greedy (\(\chref{c28}\), \(\defref{c23-map}\)): you never query a
     path on a live forest, so an LCT earns nothing. Recognise the Kruskal shape.}
  \practice{Connectivity Addicts}{CF 1738F}
    {https://codeforces.com/problemset/problem/1738/F}{CF 2400 L5}
    {Batched edge insertions followed by connectivity questions -- ``are \(u,v\) in one component?''
     with no path aggregate. That is exactly the connectivity-only row of \(\defref{c23-map}\): a
     plain or rollback DSU (\(\chref{c28}\)) answers each in \(\bigO{\alpha}\), and the tree topology
     is not even being folded. Reject the LCT; the whole cost of \code{makeroot}/\code{access} buys a
     boolean you get free.}
  \practice{Groceries in Meteor Town}{CF 1628E}
    {https://codeforces.com/problemset/problem/1628/E}{CF 3100 L6}
    {Path queries on a \emph{static} tree combined with a binary search / offline ordering -- the
     edges never change. The first question of \(\defref{c23-boundary-pick}\) (does the topology
     change?) is \emph{no}, so route it to a fixed-tree engine (HLD \(\chref{c16}\) or
     Euler+Fenwick \(\chref{c24}\)). The lesson is that ``path query on a tree'' is not by itself an
     LCT trigger -- the \emph{moving} edge is.}
  \practice{Wait for it}{CF 696E}
    {https://codeforces.com/problemset/problem/696/E}{CF 3000 L6}
    {A forest that only \emph{grows} by links while you query aggregates near the frontier -- the
     genuine link-plus-subtree shape of \(\secref{c23-beyond}\). It tempts the \(\mathit{vir}\)
     sibling of \(\coderef{c23-subtree}\); but check whether links are monotone (never cut): if so an
     Euler tour over the \emph{final} tree plus a BIT (\(\chref{c24}\)) is leaner. Decide with
     \(\defref{c23-map}\), then justify the heavier tool if you take it.}
  \practice{Pastoral Oddities}{CF 603E}
    {https://codeforces.com/problemset/problem/603/E}{CF 3000 L6}
    {Edge insertions arrive all known in advance and you maintain a component property, emitting a
     threshold edge weight -- the offline connectivity/divide-and-conquer row of \(\defref{c23-map}\).
     Because you may read every query first, a rollback DSU with D\&C on the timeline
     (\(\chref{c18}\), \(\chref{c28}\)) beats an LCT outright: \(\bigO{\alpha}\) amort per merge and a
     tenth of the code. A textbook ``the LCT is the wrong tool'' drill (\(\secref{c23-boundary}\)).}
  \practice{Surprise me!}{CF 809E}
    {https://codeforces.com/problemset/problem/809/E}{CF 3100 L6}
    {Choose a random subset of marked vertices and aggregate a path quantity over them -- the classic
     \emph{virtual tree} problem: compress the marked set to the smallest subtree spanning it and DP
     on that. This is Chapter~26's territory (\(\chref{c30}\)), \emph{not} link-cut: the base tree is
     fixed and you never rewire it. Forward-reference the virtual-tree construction at
     \(\secref{c23-boundary}\).}
  \practice{Raining season}{CF 1019E}
    {https://codeforces.com/problemset/problem/1019/E}{CF 3200 L6}
    {A fixed tree where each edge is a linear function of a query parameter and you want the
     extreme along a path -- a Li Chao / convex-hull-per-path problem. On a static tree that is HLD
     (\(\chref{c16}\)) with a Li Chao tree at each chain (\(\chref{c29}\)), never an LCT
     (\(\defref{c23-map}\)): the topology is fixed, so the extra evert machinery is dead weight while
     the hard part is the \emph{envelope}, not the path exposure.}
  \practice{Send Tree to Charlie}{CF 1254E}
    {https://codeforces.com/problemset/problem/1254/E}{CF 3300 L6}
    {A per-node state on a \emph{fixed} tree with no route you fold and no edge that moves -- a
     tree-DP / rerooting problem (\(\chref{c10}\), \(\chref{c20}\)). Fourth bullet of
     \(\secref{c23-boundary}\)'s warning box: an LCT needs a path dimension to justify itself; here
     there is none, so halving or evert buys nothing. Recognise ``a state per node, no route'' and
     hand it to DP.}
  \practice{Gregor and the Two Painters}{CF 1548E}
    {https://codeforces.com/problemset/problem/1548/E}{CF 3400 L6}
    {Not a tree problem at all: an \(n\times m\) grid whose cell \((i,j)\) has value
     \(a_i+b_j\) -- a rank-one matrix -- and you binary-search the answer, counting qualifying cells
     with a two-pointer sweep over the sorted rows and columns. It sits in this set purely as the
     sharpest rejection in \(\defref{c23-boundary-pick}\): no graph, no path, no dynamic edge. If
     you cannot immediately say ``nothing here is a forest,'' reread \(\secref{c23-boundary}\).}
\end{mixedset}


\chaptersection{The Summit -- CF 487E Tourists}

\begin{notebox}[Why this is the chapter top rope]
Every in-section problem here ran the LCT on a structure that \emph{already was} a forest. The
Summit hands you a general connected \emph{graph} and asks a path question that is not about a
single route but about \emph{some} simple route -- the exact shape \(\secref{c23-boundary}\)'s last
spec flags: ``condense first, then run the LCT.'' The biconnected (block-cut) tree is the condensation
that makes the question a path fold: once a block node carries the cheapest city inside it, the answer
is a plain path-\(\min\) on a forest -- \code{makeroot} then \code{access} (\(\secref{c23-makeroot}\)).
It is the honest summit, not the highest badge: it is the hardest problem this chapter fully derives
and hand-replays to the judge's output. \emph{One caveat the pick already told you:} the block-cut
tree here is \emph{static} (edges never change, only prices do), so a fixed-tree HLD would also serve;
we instantiate the engine this chapter built -- the LCT -- and treat the reduction, not the dynamism,
as the difficulty (\(\defref{c23-map}\)).
\end{notebox}

\solvedproblem{c23-summit}{Tourists}{CF 487E}
  {https://codeforces.com/problemset/problem/487/E}{CF 3200 L6}
  {Block-cut (biconnected) tree condensation with Tarjan edge-stack; a block node stores the minimum city price among its members and a city node its own price, so the cheapest city on SOME simple a-b route is the path-minimum on the block-cut tree; maintain it online with the chapter's evert-then-access LCT path fold, a price change re-pushing the multiset minima of the (few) blocks containing that city}

\begin{problemstatement}[CF 487E Tourists]
There are \(n\) cities (numbered \(1\) to \(n\)) joined by \(m\) bidirectional roads; city \(i\)
sells a souvenir at price \(w_i\). There are \(q\) online queries of two types:

\textbf{C a w} --- the price in city \(a\) changes to \(w\).\\
\textbf{A a b} --- a tourist travels from city \(a\) to city \(b\) along a \emph{route} (a sequence of
\emph{distinct} cities with consecutive ones joined by a road); he buys the souvenir at the cheapest
city on his route. Report the \emph{minimum possible} price he can achieve, i.e.\
\[
  \min_{\text{simple routes } x_1=a,\dots,x_k=b}\ \min_{1\le i\le k} w_{x_i},
\]
over all routes from \(a\) to \(b\) that visit no city twice (the route may be just \([a]\) when
\(a=b\)).

\textbf{Input.} Line~1: \(n,m,q\) (\(1\le n,m,q\le 10^5\)). Next \(n\) values \(w_i\)
(\(1\le w_i\le 10^9\)). Next \(m\) lines a road \(a_j,b_j\) (\(a_j\ne b_j\), no duplicate road). The
graph is guaranteed connected (a route always exists). Then \(q\) lines, each \texttt{C a w} or
\texttt{A a b}.

\textbf{Output.} For each \texttt{A} query, print the answer on its own line.

\textbf{Sample 1.} triangle \(1\!-\!2\!-\!3\!-\!1\) with prices \(1,2,3\):
\texttt{A 2 3} \(\to\) \texttt{1}; \texttt{C 1 5}; \texttt{A 2 3} \(\to\) \texttt{2}.\\
\textbf{Sample 2.} \(n{=}7\), prices \(1..7\), roads \((1,2),(2,5),(1,5),(2,3),(3,4),(2,4),\)
\((5,6),(6,7),(5,7)\); the four \texttt{A} queries answer \texttt{2, 1, 5, 3}.
\end{problemstatement}

\begin{attemptgate}[Attempt it yourself first -- \lv{6}, budget a real afternoon]
Running a graph search per \texttt{A} query and taking the cheapest city on \emph{one} found path
under-estimates the tourist's freedom: within any biconnected region he can detour. And prices
change, so nothing can be baked once.
\peek{a tourist who enters a 2-vertex-connected region can reach \emph{any} city of that region
without revisiting a city -- what single number therefore summarises a whole region?}
\nudge{Condense the graph to its block-cut tree: a bipartite tree of city-nodes and block-nodes
(maximal biconnected components). Make each block-node's value the \emph{minimum} price among its
member cities and each city-node's value its own price. Then the cheapest achievable on a route
\(a\to b\) is exactly the minimum node value on the unique \(a\)-to-\(b\) path of that tree.}
\fullhint{Build the block-cut tree once with Tarjan's edge stack (block nodes \(n+1,\dots,n+B\));
\(O(n+m)\). Run the chapter's LCT \emph{on that tree} with the path-\(\min\) fold
(\(\mathit{mn}[0]=+\infty\), \(\mathit{mn}\) = min of value and the two children): \texttt{A a b}
\(=\) \code{makeroot(a); access(b); return mn[b]}. Keep, per block, a \code{multiset} of its member
cities' prices so its value is \(*\)set.begin(). A price change \texttt{C a w} updates city-node
\(a\)'s value to \(w\) and, for each of city \(a\)'s \(\bigO{\deg}\) containing block-nodes, erases
the old price from that block's multiset, inserts \(w\), and re-pushes the block-node's value via one
\code{access}. Sum over all block sizes is \(\bigO{n+m}\), so an average update touches few nodes;
the total is \(\bigO{(n+m)+q\log n\cdot\bar{d}}\). Read/write 64-bit prices. \emph{Notice the
honesty check:} the tree itself never changes, so a fixed-tree HLD + segment tree is an equally valid
-- and by \(\defref{c23-boundary-pick}\) the leaner -- choice; the reduction to a block-cut tree is
the actual hard part.}
\end{attemptgate}

\solstep{Understanding}
A ``route'' is a \emph{simple path}, and the tourist wants the cheapest city he can arrange to pass
through. The pivotal structural fact is that a biconnected component (block) is \emph{enterable and
freely movable}: for any three distinct cities \(x,y,z\) in a 2-connected graph there is a simple
path from \(x\) to \(y\) through \(z\). So a tourist who must traverse a block from entry city
\(p\) to exit city \(q\) can also swing by \emph{any} city of that block -- in particular its
cheapest one -- without ever repeating a city. Hence a block contributes the single number
\(\min\{w_v : v\in\text{block}\}\), and only its articulation cities constrain the route.

Form the block-cut tree: bipartite, city-nodes for the original cities, block-nodes for the blocks,
a city-node joined to a block-node when the city belongs to the block. Every simple route \(a\to b\)
in the graph projects to the unique \(a\)-to-\(b\) path in the block-cut tree (the articulation
cities and blocks it must cross), and conversely the block-cut path is realisable as a graph route.
The cheapest city obtainable is therefore the minimum node value along that tree path, where a city
node is valued by \(w\) and a block node by its members' minimum. This is a path-\(\min\) on a tree
with point updates -- precisely \(\secref{c23-makeroot}\)'s fold.

\solstep{Key Observations}
\begin{observation}
\textbf{Block mobility.} Inside one biconnected component, entering at \(p\) and leaving at \(q\),
the tourist can pass through \emph{any chosen city} of the component on a simple path. Thus the
component is summarised by \emph{one} value -- the minimum member price -- and no route ever needs
to know more about it. (This collapses a whole block to a single node weight.)
\end{observation}
\begin{observation}
\textbf{Answer \(=\) path-min on the block-cut tree.} A route from \(a\) to \(b\) corresponds to the
unique \(a\)-to-\(b\) path in the block-cut tree; the minimum achievable price equals
\(\min\{\text{node value} : \text{node on that tree path}\}\) -- city nodes give their own price,
block nodes give their members' minimum. The path fold \(\mathit{mn}\) of \(\thmref{c23-path-query}\)
answers it in \(\bigO{\log n}\).
\end{observation}
\begin{observation}
\textbf{A price change touches few nodes.} Changing city \(a\)'s price to \(w\) updates city-node
\(a\), and each block-node containing \(a\) whose multiset-minimum may move -- that is
\(\bigO{\deg_{\text{BCT}}(a)}\) block nodes, each maintained by a \code{multiset} so the new value is
\(*\)begin. The block-cut \emph{topology} is fixed; only node values change, which is exactly what an
LCT point-update (access, set, pull) handles.
\end{observation}

\solstep{Think It With Me}
\begin{thinkbox}
\deadend{For each \texttt{A a b}, run BFS/Dijkstra and take the cheapest city on the found shortest path}
{Two failures. First, cost is not the number of roads, so a shortest path is the wrong object.
Second, and fatal, the tourist may take a \emph{longer} simple route to touch a cheap city a shortest
path skips. The answer is not ``min over one canonical path'' but ``min over the best simple route,''
which the block structure makes a tree-path problem. \textbf{WA}.}
\deadend{Value a block-node by only its articulation cities}
{A block's cheapest city need not be an articulation point -- it can be a leaf interior to the block,
which the tourist can still reach (Observation~1). Valuing the block by just its boundary cities
misses the whole interior and gives \textbf{WA}; the node value must be the minimum over \emph{every}
member city.}
\deadend{Rebuild the block-cut tree or DFS the tree path on every query}
{The tree has \(\bigO{n+m}\) nodes and there are \(q\le 10^5\) queries -- a per-query walk is
\(\bigO{q(n+m)}\), \tle. You need the path-\(\min\) to be polylogarithmic and the price updates to be
local, which is the LCT's job; but note the topology never changes, so HLD would also be correct
(\(\defref{c23-map}\)) -- the dynamism is in the \emph{values}, not the edges.}
\aha{Condense \emph{once} to the block-cut tree (Tarjan edge stack), give each block-node its
members' multiset-minimum and each city-node its own price, then answer \texttt{A} with the chapter's
\code{makeroot(a); access(b); mn[b]} path-\(\min\), and \texttt{C} by re-pushing the \(\bigO{\deg}\)
affected block-nodes. The hard insight is the \emph{reduction}; the data structure is the fold you
already built in \(\secref{c23-makeroot}\).}
\end{thinkbox}

\solstep{The Approach Ladder}
\begin{center}\renewcommand{\arraystretch}{1.25}
\begin{tabular}{@{}c p{0.23\linewidth} p{0.29\linewidth} l c@{}}
\toprule
\textbf{\#} & \textbf{Approach} & \textbf{Idea} & \textbf{Time} & \textbf{Verdict}\\
\midrule
1 & \lv{2}\ BFS the graph per query & min city on one shortest path & \(\bigO{q(n+m)}\) & \tle \\
2 & \lv{3}\ BFS but allow any simple path & interiors are reachable, missed & -- & \textcolor{warningred}{\ding{55}\ \textbf{WA}} \\
3 & \lv{4}\ block-cut tree, DFS path-min each query & right model, linear query & \(\bigO{q(n+m)}\) & \tle \\
4 & \lv{6}\ \textbf{block-cut tree + LCT path-min} & condense once, then fold & \(\bigO{q\log n}\) & \ac \\
\bottomrule
\end{tabular}\end{center}
Rung~4 is the Summit: rung~3 has the right \emph{model} but a linear query; replacing the per-query
tree walk with the LCT path fold (\(\thmref{c23-path-query}\)) is the whole win, and the price update
becomes a local multiset re-push.

\solstep{Solution}
Read the graph and prices. Run Tarjan's biconnected-component DFS with an edge stack: whenever
\(\mathit{low}[v]\ge \mathit{dfn}[u]\), pop the stack down to the current tree edge -- those edges'
distinct endpoints are one block, given a fresh node id \(n+\mathit{bcnt}\). Link each member city to
the block node (in the LCT as well as the tree), insert its price into the block's multiset, and set
the block node's value to the multiset minimum; set each city node's value to its own price. Answer
\texttt{A a b} by \code{makeroot(a); access(b)} and reading \(\mathit{mn}[b]\). On \texttt{C a w}, for
each block containing \(a\) erase the old price and insert \(w\), then re-push that block node's value;
finally re-push city node \(a\)'s value. Everything is a path-\(\min\) with point updates, exactly the
\(\secref{c23-makeroot}\) engine.

\solstep{Proof of Correctness}
\begin{proofbox}
\textbf{(Route \(\Rightarrow\) tree path, upper bound.)} Let \(\mathcal{P}\) be the unique
\(a\)-to-\(b\) path in the block-cut tree. Any simple graph route \(R\) from \(a\) to \(b\) crosses
exactly the blocks and articulation cities of \(\mathcal{P}\) (removing an articulation city
disconnects the graph along \(\mathcal{P}\)), so every city of \(R\) lies in some block-node or is a
city-node of \(\mathcal{P}\). Hence \(\min_{v\in R} w_v \ge \min_{x\in\mathcal{P}} \mathit{val}(x)\):
a city is either a city-node (its value is \(w\)) or inside a block whose node value is the block's
\(\min\le w\). So no route beats \(\min\mathcal{P}\). \textbf{(Tree path \(\Rightarrow\) route, lower
bound.)} Write \(\mathcal{P}=a, B_1, c_1, B_2,\dots,b\). Concatenate, for each block \(B_i\), a simple
path inside \(B_i\) from its entry articulation to its exit articulation that \emph{passes through}
\(B_i\)'s cheapest city -- realisable because a block is 2-connected (Observation~1), and these sub-paths
meet only at articulation cities, so their union is a genuine simple route. Choosing the cheapest city
of the block attaining \(\min\mathcal{P}\) shows the value \(\min\mathcal{P}\) is achievable. \textbf{(Updates.)}
The block-cut tree is topologically fixed (the graph's edges never change), so \(\mathcal{P}\) is a
function of only \(a,b\); a price change alters one city-node value and, per containing block, that
block's multiset minimum, which is exactly the set of node values on future paths that the change can
affect. The LCT \(\mathit{mn}\) fold returns \(\min\mathcal{P}\) after each update, hence every
\texttt{A} prints the true minimum. \(\Box\)
\end{proofbox}

\solstep{Implementation}
\begin{lstlisting}[caption={CF 487E Summit: Tarjan block-cut tree (edge stack) built once, then the chapter's evert-then-access LCT path-min over it. A block node stores the multiset-min of its member city prices, a city node its own price; A a b = makeroot(a);access(b);mn[b], C a w re-pushes the few blocks containing a plus city a itself. ASCII only, identifiers without underscores.}, label={lst:c23-summit}]
#include <bits/stdc++.h>
using namespace std;
typedef long long ll;
const int MAXN = 200005;                 // n cities + <= n blocks
const ll  INFL = (ll)4e18;
int n, m, q;
// ---------- LCT with path-MIN ----------
int ch[MAXN][2], fa[MAXN]; bool revTag[MAXN];
ll val[MAXN], mn[MAXN];
bool isRoot(int x){ int f=fa[x]; return f==0 || (ch[f][0]!=x && ch[f][1]!=x); }
void pull(int x){ mn[x]=val[x];
    if(ch[x][0]) mn[x]=std::min(mn[x],mn[ch[x][0]]);
    if(ch[x][1]) mn[x]=std::min(mn[x],mn[ch[x][1]]); }
void pushRev(int x){ if(!x) return; std::swap(ch[x][0],ch[x][1]); revTag[x]^=1; }
void push(int x){ if(revTag[x]){ pushRev(ch[x][0]); pushRev(ch[x][1]); revTag[x]=0; } }
void pushPath(int x){ if(!isRoot(x)) pushPath(fa[x]); push(x); }
void rot(int x){ int y=fa[x], z=fa[y], d=(ch[y][1]==x);
    if(!isRoot(y)) ch[z][ch[z][1]==y]=x; fa[x]=z;
    ch[y][d]=ch[x][d^1]; if(ch[x][d^1]) fa[ch[x][d^1]]=y;
    ch[x][d^1]=y; fa[y]=x; pull(y); pull(x); }
void splay(int x){ pushPath(x);
    while(!isRoot(x)){ int y=fa[x], z=fa[y];
        if(!isRoot(y)) rot( (ch[y][0]==x)==(ch[z][0]==y) ? y : x );
        rot(x); } }
void access(int x){ for(int last=0; x; last=x, x=fa[x]){ splay(x); ch[x][1]=last; pull(x);} }
void makeroot(int x){ access(x); pushRev(x); }
void link(int u,int v){ makeroot(u); fa[u]=v; }
void setVal(int x,ll w){ access(x); val[x]=w; pull(x); }
ll  pathMin(int u,int v){ makeroot(u); access(v); return mn[v]; }
// ---------- Block-cut tree (Tarjan edge stack) ----------
vector<pair<int,int> > g[MAXN]; int eu[MAXN], ev[MAXN];
int dfn[MAXN], low[MAXN], tmr, estk[MAXN], etop, bcnt;
ll wcity[MAXN];
multiset<ll> ms[MAXN];                     // per block node: member city prices
vector<int> ofBlock[MAXN];                 // city -> block nodes containing it
void dfs(int u,int pe){ dfn[u]=low[u]=++tmr;
    for(size_t k=0;k<g[u].size();k++){ int v=g[u][k].first, id=g[u][k].second;
        if(!dfn[v]){ estk[++etop]=id; dfs(v,id); low[u]=std::min(low[u],low[v]);
            if(low[v]>=dfn[u]){                       // peel one block
                ++bcnt; int node=n+bcnt; vector<int> verts;
                while(true){ int e=estk[etop--];
                    if(eu[e]!=0){verts.push_back(eu[e]);} // dedup below
                    verts.push_back(ev[e]);
                    if(e==id) break; }
                sort(verts.begin(),verts.end()); verts.erase(unique(verts.begin(),verts.end()),verts.end());
                for(size_t t=0;t<verts.size();t++){ int v=verts[t];
                    ms[node].insert(wcity[v]); ofBlock[v].push_back(node); link(v,node); }
                val[node]=*ms[node].begin(); pull(node);
            }
        } else if(id!=pe && dfn[v]<dfn[u]){ estk[++etop]=id; low[u]=std::min(low[u],dfn[v]); }
    } }
int main(){ scanf("%d%d%d",&n,&m,&q);
    for(int i=1;i<=n;i++){ scanf("%lld",&wcity[i]); val[i]=wcity[i]; mn[i]=wcity[i]; }
    for(int i=1;i<=m;i++){ int a,b; scanf("%d%d",&a,&b); eu[i]=a; ev[i]=b;
        g[a].push_back(make_pair(b,i)); g[b].push_back(make_pair(a,i)); }
    for(int i=1;i<=n;i++) if(!dfn[i]) dfs(i,-1);      // build BCT + LCT forest once
    for(int qi=0;qi<q;qi++){ char t; int a,b; ll w;
        scanf(" %c%d%d",&t,&a,&b);
        if(t=='A'){ printf("%lld\n", pathMin(a,b)); }
        else { w=(ll)b; ll old=wcity[a];
            for(size_t k=0;k<ofBlock[a].size();k++){ int node=ofBlock[a][k];
                multiset<ll>::iterator it=ms[node].find(old); if(it!=ms[node].end()) ms[node].erase(it);
                ms[node].insert(w); setVal(node,*ms[node].begin()); }
            wcity[a]=w; setVal(a,w); } }
    return 0; }
// The graph is condensed ONCE (topology fixed); A is a makeroot+access path-min, C is a local
// multiset re-push on the few blocks touching the city. mn[0] unused (pull reads only real children).
\end{lstlisting}

\solstep{Trace on the Official Sample}
\begin{tracebox}[Sample 2: condense, then replay the four A queries as path-minima]
Prices \(w_i=i\). The three triangles are blocks; cities \(2\) and \(5\) are the articulations.
Block nodes: \(B_a{=}8=\{1,2,5\}\) min \(=1\); \(B_b{=}9=\{2,3,4\}\) min \(=2\);
\(B_c{=}10=\{5,6,7\}\) min \(=5\). Tree edges: \(8\!-\!1,8\!-\!2,8\!-\!5,9\!-\!2,9\!-\!3,9\!-\!4,
10\!-\!5,10\!-\!6,10\!-\!7\). Node values: city \(i\mapsto i\); \(8\mapsto1,9\mapsto2,10\mapsto5\).

\textbf{A 2 3.} Tree path \(2\!-\!9\!-\!3\); values \(2,2,3\); min \(=\boxed{2}\). \(\checkmark\)

\textbf{A 6 4.} Path \(6\!-\!10\!-\!5\!-\!8\!-\!2\!-\!9\!-\!4\); values \(6,5,5,1,2,2,4\);
min \(=\boxed{1}\) (realised by detouring through city \(1\) inside block \(8\)). \(\checkmark\)

\textbf{A 6 7.} Path \(6\!-\!10\!-\!7\); values \(6,5,7\); min \(=\boxed{5}\). \(\checkmark\)

\textbf{A 3 3.} Path is the single node \(3\); value \(w_3=\boxed{3}\). \(\checkmark\) -- matching the
judge exactly. (Sample 1: the whole triangle is one block \(B\) with min \(=\min(1,2,3)=1\);
\texttt{A 2 3} walks \(2\!-\!B\!-\!3\) giving \(\min(2,1,3)=1\); after \texttt{C 1 5} the block min
becomes \(\min(5,2,3)=2\), so \texttt{A 2 3} gives \(\min(2,2,3)=2\).)
\end{tracebox}

\solstep{Complexity}
\begin{complexitybox}
Tarjan block-cut construction is \(\bigO{n+m}\) time and space; the tree has
\(\le n+m\le 2\cdot10^5\) nodes (each block costs one node, total block sizes \(\bigO{n+m}\)). Each
\texttt{A} is one \code{makeroot}/\code{access}, \(\bigO{\log N}\) amortised (\(\thmref{c23-splay-amort}\)).
Each \texttt{C a w} touches city node \(a\) plus its \(\bigO{\deg_{\text{BCT}}(a)}\) block nodes, each a
\(\bigO{\log}\) multiset edit and a \(\bigO{\log N}\) LCT \code{setVal}; summing over all price changes
that is \(\bigO{\sum_a \deg(a)}=\bigO{n+m}\) per full sweep, so on typical data the queries cost
\(\bigO{q\log N}\) and the worst single city (a hub in many blocks) is \(\bigO{\deg\log N}\) --
comfortably inside 2 seconds at \(10^5\). Prices up to \(10^9\) fit in \code{long long};
\(\mathit{mn}\) is a genuine \(\min\), and the null child is read only after checking it exists, so no
neutral-value bug can surface.
\end{complexitybox}

\begin{hardeningbox}
\begin{itemize}[nosep,leftmargin=1.4em]
  \item \textbf{A block's value is the min over \emph{all} its cities, not just articulations.} A
        leaf interior to a block is reachable (2-connectivity), so restricting the min to boundary
        cities is the classic \textbf{WA} on this problem (\(\probref{c23-summit}\)'s second dead end).
  \item \textbf{Re-push the containing blocks, then the city, on \texttt{C}.} The block multiset min
        \emph{and} the city node both changed; forgetting the city node misses routes that stop at
        \(a\) itself, forgetting a block misses the interior-min detour (\(\coderef{c23-summit}\)).
  \item \textbf{Pop the edge stack \emph{down to the current tree edge} and dedup endpoints.} A block
        is a set of edges; its vertex set is their distinct endpoints. Missing the \(\mathit{low}\ge
        \mathit{dfn}\) boundary or forgetting to dedup double-counts a city into one block's multiset.
  \item \textbf{Keep \(\mathit{mn}\) a pure \(\min\) and guard null children.} Reading
        \(\mathit{mn}[0]\) as \(0\) (not \(+\infty\)) would sink every path-min to \(0\); the fold is
        a min, so the neutral is \(+\infty\) and \code{pull} must only fold existing children.
  \item \textbf{The topology is static -- but that does not excuse a wrong engine claim.} Because no
        edge is ever added, an honest reader notes HLD + a segment tree also passes
        (\(\defref{c23-boundary-pick}\)); the LCT is used to reuse \(\secref{c23-makeroot}\), not
        because the forest moves. Say so under pressure rather than over-claiming dynamism.
\end{itemize}
\end{hardeningbox}

\begin{recallbox}[Chapter 25 in ten seconds]
A splay shuffles the touched key to the root in \(\bigO{\log n}\) amortised (\(\thmref{c23-splay-amort}\));
a link-cut tree chains splays along preferred paths so \code{access} exposes a root-to-node path in
one loop (\(\defref{c23-access}\), \(\coderef{c23-lct-core}\)). Rooted, the two-access climb reads an
LCA (\(\thmref{c23-access-lca}\), CF DYNALCA \(\coderef{c23-dynalca}\)); with \code{makeroot} and a
reverse tag, any path fold drops out of \code{access} (\(\thmref{c23-path-query}\),
\(\coderef{c23-evert}\)); a virtual-child accumulator extends it to subtree aggregates
(\(\thmref{c23-vsum}\), \(\coderef{c23-subtree}\)). Fixed trees want HLD/Euler, offline connectivity
wants DSU, all-pairs counting wants centroid (\(\defref{c23-boundary-pick}\)). Summit CF 487E
(\(\probref{c23-summit}\)) condenses a graph to its block-cut tree and reads a path-min.
\selftest{
  \item Why is \(\mathit{isRoot}(x)\) not simply \(\mathit{fa}[x]=0\)? \hfill\lv{4}
  \item State the two-access LCA procedure and why the last splayed node is the ancestor. \hfill\lv{5}
  \item Why must \code{pushPath} run before the zig-zig loop in \code{splay}? \hfill\lv{5}
  \item What does \(\mathit{vir}[x]\) hold, and why does the path fold omit it while the subtree fold
        includes it? \hfill\lv{5}
  \item In the Summit, what value does a block node carry and why is it not just its articulations? \hfill\lv{6}
}
\end{recallbox}


\chaptersection{Tool Kit Inventory}

The machine assembled through Chapter~25. Earlier chapters' rows are carried forward; this
chapter's are appended. Pointers only.

\begin{center}
\small
\renewcommand{\arraystretch}{1.25}
\begin{tabular}{@{}p{0.23\linewidth} l c p{0.22\linewidth} >{\raggedright\arraybackslash\scriptsize}p{0.18\linewidth}@{}}
\toprule
\textbf{Tool} & \textbf{Since} & \textbf{Lv} & \textbf{Trigger phrase} & \textbf{Where} \\
\midrule
\multicolumn{5}{l}{\itshape Carried forward} \\
Euler tour / subtree-as-interval & 5 & \lv{3} & ``a fixed tree's subtree is one contiguous range'' & \chref{c24} \\
LCA + binary lifting & 16 & \lv{4} & ``the static-tree ancestor the rooted LCA replaces'' & \chref{c11} \\
Heavy-light decomposition & 22 & \lv{5} & ``a path on a tree that never changes'' & \chref{c16} \\
DSU / rollback DSU & 17 & \lv{4} & ``connectivity only, offline'' & \chref{c28} \\
Centroid decomposition & 24 & \lv{5} & ``counting over ALL paths of a fixed tree'' & \chref{c18} \\
\midrule
\multicolumn{5}{l}{\itshape Chapter 25 -- new this chapter} \\
Splay tree (amortised) & 25.1 & \lv{4} & ``shuffle the touched key up; order statistics'' & \defref{c23-splay}, \coderef{c23-splay-rot} \\
Top-down splay & 25.1 & \lv{4} & ``one downward pass, no parent pointers'' & \coderef{c23-topdown} \\
Preferred paths + access & 25.2 & \lv{4} & ``expose root-to-x on a moving forest in one loop'' & \defref{c23-access}, \coderef{c23-lct-core} \\
Two-access LCA & 25.2 & \lv{5} & ``LCA / connectivity on a growing rooted forest'' & \thmref{c23-access-lca}, \coderef{c23-dynalca} \\
Evert + path fold & 25.3 & \lv{5} & ``sum/min over the current u--v path, rerootable'' & \thmref{c23-path-query}, \coderef{c23-evert} \\
Virtual-child subtree fold & 25.4 & \lv{5} & ``a subtree aggregate on a moving tree'' & \thmref{c23-vsum}, \coderef{c23-subtree} \\
The dynamic-tree engine pick & 25.5 & \lv{5} & ``is the topology even moving, and do I need a path value?'' & \defref{c23-boundary-pick}, \defref{c23-map} \\
Link-cut kit (one listing) & 25 & \lv{5} & ``all the skeletons in one place'' & \coderef{c23-kit} \\
Summit: block-cut + LCT path-min & Summit & \lv{6} & ``cheapest city on some simple route, prices change'' & \coderef{c23-summit} \\
\bottomrule
\end{tabular}
\end{center}

\chaptersection{Related Topics}

\begin{itemize}[nosep, leftmargin=1.4em]
  \item \chref{c16} (Ch.~22, Heavy-Light Decomposition) --- contrast and promise \textbf{discharged}:
        HLD folds a path on a \emph{fixed} tree in \(\bigO{\log^2 n}\); \(\secref{c23-makeroot}\) folds a
        path on a \emph{moving} forest in \(\bigO{\log n}\) amortised. \(\secref{c16-boundary}\) handed
        the moving-topology case forward to this chapter; \(\defref{c23-boundary-pick}\) makes the split
        explicit -- the tell is whether an edge ever appears or vanishes.
  \item \chref{c11} (Ch.~16, LCA) --- backward/contrast: binary lifting answers a static LCA in
        \(\bigO{\log n}\) but cannot survive a link or cut; the two-access LCA
        (\(\thmref{c23-access-lca}\), \(\coderef{c23-dynalca}\)) is its dynamic sibling, recomputing the
        ancestor as the rooted forest rewires.
  \item \chref{c18} (Ch.~24, Centroid Decomposition) --- contrast: centroid counts over \emph{all} paths
        of a fixed tree; the LCT folds \emph{one} named path of a changing tree. Both die if you mix the
        scopes -- centroid needs a static topology, LCT needs a moving one (\(\defref{c23-map}\)).
        \textbf{This is the Previous chapter.}
  \item \chref{c24} (Ch.~5, Euler Tour / tree representations) --- backward: the Euler tour is the
        \emph{static} way a subtree becomes an interval; an Euler-tour tree keeps that tour in a balanced
        BST and is the simpler dynamic-forest tool when you only need connectivity, not a path fold
        (\(\secref{c23-beyond}\)).
  \item \chref{c28} (Ch.~17, DSU / Sack) --- forward/contrast: a rollback DSU answers offline
        connectivity in \(\bigO{\alpha}\) with a fraction of the code; reach for it before an LCT whenever
        the edits are all known in advance (\(\defref{c23-map}\)).
  \item \chref{c29} (Ch.~23, Li Chao / Segment-Tree Beats) --- forward/contrast: a path whose aggregate is
        a line envelope evaluated at a variable is Li Chao hosted on HLD chains (CF 1019E in the mixed
        set), not a monoid fold on an LCT.
  \item \chref{c10} (Ch.~13, Tree DP) --- contrast: a state \emph{per node} with no route dimension is DP
        / rerooting (\(\chref{c20}\)); the LCT needs a path (or subtree) you fold, so halving/evert buys
        nothing on pure DP specs (\(\secref{c23-boundary}\)).
  \item \chref{c30} (Ch.~26, Virtual Trees) --- forward: the mixed-set CF 809E hands a random-marked-subset
        aggregate to the virtual tree, the next chapter's engine -- compressing a fixed tree to the nodes
        that matter rather than making it dynamic. \textbf{This is the Next chapter.}
  \item \chref{c21} (Ch.~27, Mo's Algorithm) --- forward: an offline batch of node-set queries with no
        inverse is Mo-on-trees, complementing the LCT's online, invertible path folds
        (\(\secref{c23-boundary}\)).
\end{itemize}

\begin{lastword}[Last Word --- Chapter 25]
\emph{``A splay tree remembers what you touched last and charges you for it in the right order; a
link-cut tree chains those splays along preferred paths and re-preferences them with a single access,
so a forest that keeps rewiring its edges can still answer a path -- and, with a virtual-child ledger,
a subtree -- in amortised \(\bigO{\log n}\).''}

\medskip
Chapter~22 folded paths only on a tree that stood still; this chapter dropped that assumption. Section
25.1 built the splay and its potential-function bound (\(\thmref{c23-splay-amort}\)); 25.2 turned
splays into a forest with \code{access} and read an LCA two-accesses at a time
(\(\thmref{c23-access-lca}\)); 25.3 added \code{makeroot} and a lazy reversal so any path fold drops out
of \code{access} (\(\thmref{c23-path-query}\)); 25.4 paid the honest price of a \emph{subtree} on a
moving tree with the \(\mathit{vir}\) ledger (\(\thmref{c23-vsum}\)); and the Summit CF 487E
(\(\probref{c23-summit}\)) showed that sometimes the tree you need is one you must \emph{build} --
condensing a graph to its block-cut tree before the fold even applies.

\medskip
Carry three reflexes forward:
\begin{enumerate}[nosep, leftmargin=1.6em]
  \item \emph{Ask whether the topology actually moves.} If no edge is ever linked or cut, you wanted
        HLD or Euler+Fenwick, not an LCT (\(\defref{c23-boundary-pick}\)); the moving edge is the whole
        licence to pay splay constants.
  \item \emph{Access is the one primitive.} \code{findroot}, \code{makeroot}, \code{link}, \code{cut}
        and every path or subtree fold are the same root-to-node exposure wearing different hats
        (\(\defref{c23-access}\), \(\coderef{c23-kit}\)).
  \item \emph{Push your lazy tags before you rotate.} A splay that rotates under an un-pushed reverse
        tag reattaches mirrored children -- the deepest, rarest WA in the structure
        (\(\defref{c23-makeroot}\), \(\coderef{c23-evert}\)).
\end{enumerate}

\medskip
And three habits, drilled until boring: keep \(\mathit{isRoot}\) a splay test, never \code{fa==0}; keep
the path fold \(\mathit{sum}\) and the subtree fold \(\mathit{sub}\) as disciplined siblings so a query
can never grab the wrong one; and when a ``dynamic tree'' problem hands you a \emph{graph}, reach for a
condensation (block-cut or spanning forest) before you reach for \code{makeroot}. You now own the one
tree engine that answers while the tree itself changes -- the dynamic companion to every static
structure in Part~V.

\medskip
\textbf{Next:} \chref{c30} (Virtual Trees) turns the difficulty inside-out. Instead of a
\emph{moving} tree with all of \(n\) nodes, you get a \emph{fixed} tree and a small marked subset per
query, and the win is to compress the tree down to just the marked nodes and their LCAs -- an
\(\bigO{k\log k}\) skeleton -- then run a DP on it. Where this chapter made the structure move, the
next one makes it \emph{shrink}, and it cashes in the LCA machinery of Chapter~16 that this chapter
merely borrowed.
\end{lastword}
